\documentclass[11pt]{article}

\usepackage[utf8]{inputenc}
\usepackage[T1]{fontenc}
\usepackage{amsmath, amssymb, amsthm}
\usepackage{booktabs}
\usepackage[margin=1in]{geometry}
\usepackage{hyperref}

\newtheorem{theorem}{Theorem}

\title{A Short Example Paper}
\author{Your Name \\ \small Your Institution}
\date{\today}

\begin{document}

\maketitle

\begin{abstract}
This example shows the pieces most papers need: sections, equations,
a theorem, a table, cross-references, and a bibliography. Edit any part
of it, compile, and watch the PDF update.
\end{abstract}

\section{Introduction}
\label{sec:intro}

Writing in \LaTeX{} separates content from layout. Cross-references such
as Section~\ref{sec:method} and Equation~\eqref{eq:energy} update by
themselves, and citations come from a \texttt{.bib} file~\cite{knuth1984}.

\section{Method}
\label{sec:method}

Consider a particle of mass $m$ moving with velocity $v$. Its kinetic
energy is
\begin{equation}
  E_k = \frac{1}{2} m v^2 .
  \label{eq:energy}
\end{equation}

\begin{theorem}
For every real $x$, $x^2 \geq 0$.
\end{theorem}

\begin{proof}
If $x \geq 0$ the product of two non-negative numbers is non-negative;
otherwise $x^2 = (-x)^2$ with $-x > 0$.
\end{proof}

\section{Results}

Table~\ref{tab:results} lists a few made-up measurements.

\begin{table}[h]
  \centering
  \begin{tabular}{lrr}
    \toprule
    Trial & Mass (kg) & Energy (J) \\
    \midrule
    A & 1.0 & 2.0 \\
    B & 2.5 & 7.8 \\
    C & 4.0 & 18.0 \\
    \bottomrule
  \end{tabular}
  \caption{Example measurements.}
  \label{tab:results}
\end{table}

\section{Conclusion}

Replace this text with your own work.

\bibliographystyle{plain}
\bibliography{references}

\end{document}
